\scanpage{7}
UNIVERSITE\\
DES SCIENCES ET TECHNIQUES\\
DU LANGUEDOC\\
\texttt{------------}\\
Service de la Scolarité\\
MA/GJ\\
Diffusion : Liste A\\
N\textdegree\ : 5983

Montpellier, le 12 Novembre 1980

AVIS DE SOUTENANCE D'UNE THESE DE DOCTEUR INGENIEUR

SCIENCES ALIMENTAIRES Option GENIE ET TECHNOLOGIE ALIMENTAIRES -

FERMENTATIONS

\texttt{----------------------}

\uline{Nom et prénom du candidat} : COMTE Francis

\uline{Titre de la thèse} : Etude des isothermes d'adsorption de composés phénoliques.

\uline{Date, heure et lieu de soutenance prévus} : 18 NOVEMBRE 1980 à 15 Heures 30 mn\\
Salle de cours n\textdegree\ 9 - U.S.T.L.

\uline{Nom du Directeur de thèse} : Monsieur le Professeur J. CROUZET

\uline{Désignation du laboratoire où la thèse est préparée} : Biochimie Appliquée

\uline{Résumé de la thèse} :

Le tracé d'isothermes d'adsorption sur silice de composés phénoliques (acides benzoïque et salycilique) étudiés en qualité de modèles d'anthocyanes a été obtenu par des méthodes dynamiques mettant en oeuvre la chromatographie liquide haute performance.

Parmi les différentes méthodes utilisées par élution (profil de pic, maximum de pics) ou par analyse frontale (équilibre, profil de désorption) la technique du profil de désorption est la plus rapide mais sujette à cautions. La méthode la plus fiable, quoique la plus longue, est la méthode à l'équilibre, c'est d'ailleurs la seule méthode qui autorise le tracé des isothermes individuelles par adsorption d'un mélange.

Par modélisation nous avons pu montrer que le meilleur accord était obtenu avec le modèle de SCOTT et KUCERA qui décrivent une adsorption en double couche.
En dépit des difficultés qui se sont liées à l'intervention de phénomènes propres à la chromatographie liquide et qui contribuent à la déformation des pics il nous a été possible de montrer que la capacité d'adsorption de la silice était supérieure pour l'acide ben-\\
zoïque. Ceci est confirmé par l'étude de l'adsorption mixte sans que l'on puisse parler d'adsorption spécifique ou de sélectivité de l'absorbant.

\uline{Durée approximative des recherches ayant abouti à la thèse} : 3 ans

\uline{Composition du jury proposé} :\\
MM. ROUZAUD J.\\
\hspace*{2em}CROUZET J.\\
\hspace*{2em}GOMA G.\\
\hspace*{2em}COMBE P.

P. le Secrétaire Général,\\
L'Attachée d'Administration Universitaire,

\uline{\unsure{P. Gentil|A. Gentil}}\\
\uline{A. GENTIL}

\figure{\begin{tikzpicture}
\draw[gray,densely dotted] (0.1,0.7)..controls(-0.3,0.3)and(-0.2,-0.1)..(0.05,-0.35);
\draw[gray] (0.13,0.65)--(0.17,0.25);
\draw[gray,densely dotted] (0.8,0.5)--(0.9,0.3);
\end{tikzpicture}}

\mspage{6}
